%===============================================================================
\section{Introduction}\label{sec:intro}
%===============================================================================

\IEEEPARstart{E}{lectricity} markets publish a large and growing amount of operational data. Disclosure rules list, field by
field, what must be released and when: forecasts and requirements before clearing, prices and their components right after
clearing, and actual operating quantities on the following day~\cite{nea2024disclosure}. Other quantities---the output or
the offers of an individual generating unit, the hourly flow on a specific line, or the actual load before it is
disclosed---are protected, either permanently or until a later release date. Data classification standards ask operators
to grade each field and to consider how aggregation and processing change its grade~\cite{gbt43697}, but they do not say how.

The difficulty is that the inference risk of a field is not a property of the field alone. Power data come from one physical
and economic system, and the information a field adds depends on what else is known. A unit's hourly output says little about
its private offer, and neither does the price at its bus; together, whenever the unit is marginal, they reveal the offer
(Section~\ref{sec:res-mech}). This is the inference and aggregation problem of database
security~\cite{lunt1989aggregation,farkas2002inference}, and it is not hypothetical for power data: public prices alone can
expose grid topology~\cite{kekatos2014topology}. Recent work on data sensitivity assessment has recognized that the risk of a
field depends on its inference relations to sensitive data and propagates risk along an inference graph~\cite{hu2026ignn}.
What is still missing is a field-level risk that is \emph{defined} by what an adversary can infer from a set of fields---so
that it can be checked by an attack experiment rather than calibrated to expert scores---and that captures inference
requiring several fields jointly.

We start from the gain $\Delta_i(T)$ that releasing field $i$ brings over a background $T$ of fields the adversary already
holds, measured by the out-of-sample accuracy of retrained attackers. A field therefore has a whole distribution of
contextual gains, and grading it means summarizing that distribution. Attribution methods such as Shapley values and SAGE
average the gains over backgrounds, which answers how much a field contributes typically; a security audit asks whether
\emph{some} admissible background makes the field dangerous. We therefore take the maximum---the marginal contribution
index~\cite{catav2021mci}, in line with worst-case leakage measures in quantitative information
flow~\cite{smith2009qif,alvim2012gleakage}---and add the two conditions a disclosure audit needs: the data already public and a
bound on the adversary's background. The resulting budgeted marginal risk $\Mik$, its single-field part $M^{(0)}$, the
\emph{background amplification} $\Gamma^{(K)}=\Mik-M^{(0)}$ and the backgrounds that attain it form the field-level output of
the audit.

Computing this profile requires the inferability of every small field set---7,175 sets per target for 35 candidate fields at
$K=2$---and it must be recomputed whenever the public baseline or the target definitions change. Retraining an attacker family
for each set takes hours to days per table, and general-purpose predictors that avoid training still process the training
data anew for every set. We instead amortize over the sets of one data set: a network pre-trained once under random
visibility masks supplies a representation shared by all sets, and each query only solves a closed-form readout on the
features of its fields. The estimate is used to scan; a handful of backgrounds per field is retrained to certify the result,
so that the flags do not rest on the estimate. Pre-training on the sensitive targets alone turns out to collapse the
representation and to add little over random features; our estimator therefore combines column reconstruction with the
features of an untrained network.

The contributions of this paper are as follows.
\begin{itemize}
\item \emph{Background-aware field risk.} From set-level inferability we define the budgeted marginal risk $\Mik$ and the
background amplification $\Gamma^{(K)}$, relate them to the marginal contribution index, give a closed form of the
amplification at $K=1$, and prove that $\Mik$ is the minimal threshold-free safety margin and that its critical fields are
exactly the members of small minimal unsafe sets (Section~\ref{sec:mk}).
\item \emph{An amortized attacker with scan--certify.} Masked pre-training with column reconstruction and a per-set
closed-form readout estimate $\Vy$ for every small set in tens of milliseconds; an error analysis explains why field risk is
estimated more accurately than set values, and certification by retraining yields lower bounds independent of the estimate.
The estimator is more accurate than linearized, fine-tuned and surrogate alternatives and certifies as much of the risk as a
tabular foundation model at a small fraction of its cost (Sections~\ref{sec:model} and~\ref{sec:res-est}).
\item \emph{Mechanism and real-data evidence.} A controlled market case with a known inference mechanism, and four data sets
with rule-based field roles (RTS-GMLC, the Australian NEM with actual unit outputs, PJM and CAISO) with $1.2\times10^5$
exhaustively retrained field sets as ground truth, show that background amplification is common, physically interpretable and
invisible to single-field and attribution scores (Sections~\ref{sec:res-mech}--\ref{sec:res-robust}); withholding is
discussed as an application (Section~\ref{sec:application}).
\end{itemize}

\begin{figure*}[!t]
\centering
\includegraphics[width=\textwidth]{fig1_framework.pdf}
\caption{Audit framework. Disclosure rules fix field roles. Set-level inferability $V_y(S)$ is defined by dedicated
retraining and scanned by an amortized attacker; the scan proposes, for every field, backgrounds that are certified by
retraining. The field-level profile---single-field risk, background amplification and top backgrounds---supports grading,
withholding and re-grading.}
\label{fig:framework}
\end{figure*}
